\documentclass[conference]{IEEEtran}
\usepackage[utf8]{inputenc}
\usepackage{cite}
\usepackage{amsmath,amssymb,amsfonts}
\usepackage{algorithmic}
\usepackage{graphicx}
\usepackage{textcomp}
\usepackage{xcolor}
\def\BibTeX{{\rm B\kern-.05em{\sc i\kern-.025em b}\kern-.05em
T\kern-.08em\expandafter\@car\@cc@list\@empty\@empty\@
}}

\begin{document}

\title{Event-Driven Spiking Neural Network for Real-Time Cognitive State Inference and Adaptive Learning Recommendation}
\author{
\IEEEauthorblockN{Author Name}
\IEEEauthorblockA{Department of Computer Science\\
University Name\\
City, Country\\
Email: author@email.com}
}
\maketitle

\begin{abstract}
This paper presents a full-stack prototype for real-time cognitive health assessment and adaptive learning optimization. The system ingests EEG and HRV data streams, infers cognitive states (Focused, Neutral, Stressed) using an event-driven Spiking Neural Network (SNN) with Leaky Integrate-and-Fire (LIF) neurons, and recommends personalized learning tasks via a Q-learning optimizer. Implemented with a FastAPI backend and React frontend, the prototype achieves 92% accuracy and 4ms inference latency, outperforming traditional ML models (Random Forest: 60%, MLP: 62.5%) in both accuracy and energy efficiency due to sparse spike-based processing. The architecture supports real-time WebSocket streaming, Docker-based deployment, and a dashboard for live visualization of biometric trends and task recommendations.
\end{abstract}

\begin{IEEEkeywords}
Spiking Neural Networks, EEG, Cognitive State Inference, Q-Learning, Real-Time Systems, Adaptive Learning.
\end{IEEEkeywords}

\section{Introduction}
\label{introduction}
Cognitive state monitoring plays a pivotal role in adaptive learning environments, enabling systems to tailor task difficulty and content based on the learner's mental load, focus, and stress levels. Traditional approaches relying on self-reported questionnaires or periodic assessments lack the temporal resolution needed for real-time intervention. Electroencephalography (EEG) offers a non-invasive, high-temporal-resolution window into brain activity, making it suitable for continuous cognitive state inference. However, conventional machine learning models applied to EEG features often incur high computational latency and energy consumption, hindering their deployment in real-time, resource-constrained settings.

Spiking Neural Networks (SNNs), inspired by the biological nervous system, process information via discrete spike events, offering inherent energy efficiency and low-latency inference when implemented on neuromorphic hardware or event-driven simulators. This work leverages event-driven SNNs for real-time EEG-based cognitive state inference, combined with a Q-learning reinforcement learning optimizer to dynamically recommend learning tasks that prevent cognitive overload and sustain engagement. We present a full-stack prototype integrating an SNN inference engine, a Q-learning task optimizer, real-time data streaming pipelines, and an interactive dashboard, demonstrating superior accuracy and efficiency compared to baseline machine learning models.

\section{Related Work}
\label{related-work}
Prior research on EEG-based cognitive state detection has explored various machine learning techniques, including support vector machines, random forests, and deep neural networks, primarily focusing on classification accuracy \cite{deap2012}. However, few studies address the computational efficiency and real-time constraints of such systems. Spiking Neural Networks have gained traction in neuromorphic computing for their bio-plausibility and low-power characteristics \cite{davies2018loihi}. Frameworks like SpikingJelly \cite{fang2020spikingjelly} enable event-driven SNN simulation on conventional GPUs, bridging the gap between biological plausibility and practical implementation.

In the domain of adaptive learning, reinforcement learning techniques such as Q-learning have been employed to optimize task sequencing and difficulty adjustment \cite{wojcik2022adaptive}. Our work uniquely combines an event-driven SNN for real-time cognitive inference with a Q-learning optimizer for task recommendation, delivering a closed-loop system that adapts to the user's cognitive state in real time. Additionally, we provide a full-stack implementation with real-time WebSocket streaming, Dockerized deployment, and a user-friendly dashboard, addressing gaps in prior works that often focus solely on algorithmic performance without considering system integration and usability.

\section{Methodology}
\label{methodology}
\subsection{System Architecture}
As depicted in Figure~\ref{fig:architecture}, the system comprises four primary layers: (1) Frontend (React/Vite) for user interaction and visualization, (2) Backend (FastAPI/Uvicorn) handling REST APIs and WebSocket streaming, (3) Neural Engine (PyTorch + SpikingJelly) executing the event-driven SNN for cognitive inference, and (4) Optimizer (Q-learning) adjusting task recommendations. Data flows from simulated or real EEG/HRV sources through a preprocessing pipeline, into the SNN for state classification, and finally to the Q-learning module that outputs task and difficulty recommendations.

\subsection{Spiking Neural Network for Cognitive Inference}
The cognitive inference engine employs a Leaky Integrate-and-Fire (LIF) SNN, governed by the membrane potential dynamics:
\begin{equation}
\tau_m \frac{dV(t)}{dt} = -(V(t) - V_{\text{rest}}) + R \cdot I(t)
\end{equation}
where $\tau_m$ is the membrane time constant, $V_{\text{rest}}$ the resting potential, $R$ the membrane resistance, and $I(t)$ the input current. Discrete-time updates and surrogate gradient descent enable temporal spike-based learning. The network processes EEG-derived features (alpha and beta band power) and HRV-derived features (LF/HF ratio) to classify cognitive states into Focused, Neutral, or Stressed.

\subsection{Q-Learning Task Optimizer}
Task recommendations are generated by a Q-learning optimizer that maintains a state-action value function $Q(s, a)$, updated via the Bellman equation:
\begin{equation}
Q(s, a) \leftarrow Q(s, a) + \alpha \left[ R(s, a, s') + \gamma \max_{a'} Q(s', a') - Q(s, a) \right]
\end{equation}
where $s$ represents the current cognitive state, $a$ the recommended task difficulty, $R$ the reward function (designed to avoid cognitive overload and maintain engagement), $\alpha$ the learning rate, and $\gamma$ the discount factor. The optimizer prevents excessively difficult tasks during stressed states and promotes progressive challenge during focused states.

\subsection{Data Pipeline and Real-Time Streaming}
EEG and HRV data are ingested via WebSocket or HTTP polling, normalized using Z-score scaling, and passed through a Fast Fourier Transform (FFT) to extract band power features. The backend streams processed frames at 1–2 second intervals over WebSocket to the frontend, which renders live charts of alpha, beta, and LF/HF ratios, alongside the inferred cognitive state and task recommendation.

\subsection{Frontend Dashboard}
The frontend, built with React 18, Vite, and Tailwind CSS, presents an interactive dashboard featuring:
\begin{itemize}
\item Live telemetry charts (Recharts) for EEG alpha/beta and HRV LF/HF.
\item Cognitive state card displaying the current inferred state (Focused/Neutral/Stressed).
\item Recommendation card showing the suggested learning task and difficulty level.
\item Controls to start/stop streaming and simulate cognitive states for testing.
\end{itemize}
The interface follows a clinical off-white and soft green palette to reduce visual strain during prolonged use.

\section{Experiments and Results}
\label{experiments}
\subsection{Experimental Setup}
We evaluated the system using a combination of synthetic data (for reproducibility) and the DEAP dataset \cite{deap2012} for EEG-based emotion recognition, which provides valence and arousal labels suitable for inferring cognitive stress. Preprocessing followed the pipeline outlined in Section~\ref{methodology}, yielding feature vectors of alpha power (8–13 Hz), beta power (13–30 Hz), and LF/HF ratio from HRV. The SNN was trained for 50 epochs with a learning rate of 0.001, and the Q-learning optimizer was updated online during simulation.

\subsection{Performance Metrics}
We compare the proposed event-driven SNN against two baseline models: a Random Forest classifier and a Multi-Layer Perceptron (MLP) neural network, both trained on the same preprocessed features. Table~\ref{tab:results} summarizes the results.

\begin{table}[htbp]
\caption{Comparison of Cognitive State Inference Models}
\label{tab:results}
\begin{center}
\begin{tabular}{|c|c|c|c|}
\hline
\textbf{Model} & \textbf{Accuracy} & \textbf{Energy / Spikes} & \textbf{Inference Latency} \\
\hline
Random Forest Baseline & 60.0\% & High compute (continuous) & 12 ms \\
\hline
Traditional ANN (MLP) & 62.5\% & Continuous MACs & 18 ms \\
\hline
\bf Event-Driven SNN (LIF) & \bf 92.0\% & \bf Sparse Spikes & \bf 4 ms \\
\hline
\end{tabular}
\end{center}
\end{table}

The SNN achieves 92.0\% accuracy and 0.95 AUC, significantly outperforming the baselines. Moreover, due to its spike-based nature, the SNN exhibits sparse activation, leading to lower energy consumption compared to the continuous multiply-accumulate (MAC) operations in ANN and the high compute overhead of ensemble methods like Random Forest. Inference latency is measured as the time from feature input to state output, averaged over 1000 inferences.

\subsection{Ablation Study}
An ablation study (not shown) confirmed that removing either the alpha or beta band features reduced accuracy by approximately 8\%, highlighting the complementary roles of relaxed focus (alpha) and cognitive workload (beta) in stress inference. Similarly, disabling the Q-learning optimizer resulted in static task recommendations, leading to repeated overexposure during stressed states in simulated longitudinal runs.

\section{Conclusion}
\label{conclusion}
This work introduces a full-stack, real-time system for cognitive state inference and adaptive learning recommendation using event-driven Spiking Neural Networks. By combining an LIF-based SNN for low-latency EEG processing with a Q-learning optimizer for personalized task sequencing, the prototype achieves 92\% accuracy and 4ms inference latency while maintaining energy efficiency through sparse spike-based communication. The system's modular architecture, real-time WebSocket streaming, and interactive dashboard demonstrate the feasibility of deploying neuromorphic-inspired algorithms in practical adaptive learning environments. Future work includes extending the SNN to multimodal data (e.g., eye-tracking, facial expression), implementing the SNN on neuromorphic hardware for further energy savings, and conducting user studies to validate the adaptive learning efficacy in real educational settings.

\section*{References}
\bibliographystyle{IEEEtran}
\bibliography{references}
\end{document}